\section{Impact Accounting: Notation and Motivation}\label{sec:formalism}
% provenance: jrn_01KXT490ZD4444AWMRG5296GJA (original radius definitions).
% provenance: jrn_01KXTVC4CT850YTVE5EM3RJDGS (PI revision directive: state-dependent flexibility, 3-part persistence, optimization BR_phys).
% provenance: dec_01KXT4C6ZPYD1Z92DH6CN1KKM3 fixes the primary physical metric.
% M10(b): compressed from a claimed formal contribution to notation and motivation only.

This section fixes notation. It is not offered as a contribution: as
Section~\ref{sec:separation} makes explicit, the accounting below is what one obtains by adding
resource-by-scope product nodes, a time index, and attached state and policy to a classical attack
graph. We need a vocabulary for ``how much, for how long, under which grid conditions'', because
the experiments in Section~\ref{sec:prelim} move those three quantities differently, and a single
scalar ``blast radius'' cannot express that. We keep the $\mathrm{BR}$ prefix for continuity with
prior drafts and call each quantity a dimension.

\paragraph{Substrate.} A static privilege graph $G=(V,E)$ has principals and resources as nodes and
access relations as edges. A compromise $x \subseteq V$ is a set of compromised credentials, and a
protection policy $\pi=(L,\mathrm{Att},A)$ encodes a credential-lifetime function $L$, an
attestation-gated renewal predicate $\mathrm{Att}$, and an authorization scope map $A$. The device
and grid operating state at time $t$ is $u_t \in U$ and $F$ is the feeder power-flow and protection
model. Neither $u_t$ nor $F$ appears in $G$, which is the whole point.

\paragraph{The four dimensions.}
\emph{Cyber reach} $\BRreach = \Reach(x,\pi)$ is the set of endpoints to which the compromised
authority can submit an authenticated request, the least fixed point of the $\pi$-enabled access
relation from $x$. Reaching a server is not controlling a resource, so we keep it separate from
everything below.

\emph{Commandable flexibility} $\BRflex$ is state dependent rather than nameplate. For DER $d$,
$\mathcal{C}_d(t;u_t,A_\pi(d))$ is the set of set-points an adversary holding scope $A_\pi(d)$ can
command given irradiance, state of charge, present set-point, and inverter apparent-power and ramp
limits. Aggregating over reachable DERs gives achievable intervals
$[\Delta P^{-},\Delta P^{+}]$ and $[\Delta Q^{-},\Delta Q^{+}]$; where a scalar is needed we use
$\BRflex = w_P(\Delta P^{+}{-}\Delta P^{-}) + w_Q(\Delta Q^{+}{-}\Delta Q^{-})$ and state the
weights. A read-only scope yields the singleton current set-point, hence zero flexibility.

\emph{Retained authority} separates three durations that certificate lifetime alone conflates:
$T_{\mathrm{cred}}$, until the stolen credential is refused for a \emph{new} session;
$T_{\mathrm{sess}}$, until sessions already authenticated with it can no longer issue commands; and
$T_{\mathrm{cmd}}$, until queued, scheduled, or latched controls stop affecting the DER. Then
\begin{equation}
\BRauth = \max\big(T_{\mathrm{cred}}, T_{\mathrm{sess}}, T_{\mathrm{cmd}}\big),
\qquad
\BRexp = \int \BRflex(t)\,\mathbf{1}[\text{authority effective at } t]\, dt,
\end{equation}
with $\BRexp$ in kW$\cdot$h or kVAr$\cdot$h. This distinction is the design's load-bearing one:
expiry bounds $T_{\mathrm{cred}}$ only, so shortening lifetime without enforcing $T_{\mathrm{sess}}$
and $T_{\mathrm{cmd}}$ does not bound exposure.

\emph{Physical consequence} $\BRphys = \max_{a\in\mathcal{A}} J(F,u,a)$ is the worst outcome
reachable within the adversary's authority, where $\mathcal{A}$ is constrained to authorized
functions and device limits. Our primary objective $J_V$ is the voltage-violation area, the
node-time outside the ANSI~C84.1 band weighted by magnitude. Secondary metrics are maximum
per-unit deviation, violating-node count, protection trips, regulator tap operations, energy not
served, unnecessary curtailment, and recovery time. We report both an optimizing attacker and an
operationally realistic one, so no result is only an omniscient upper bound.

Table~\ref{tab:dimensions} gives the types and units. The types differ --- a set, a power region, a
set of durations, and a physics functional --- which is why we do not collapse them into one
number.

\begin{table}[t]
\centering
\caption{The four impact dimensions. $\BRphys$ carries two units depending on whether it is
evaluated at one operating point or accumulated over a horizon; both appear in
Section~\ref{sec:prelim} and each table states which.}
\label{tab:dimensions}
\footnotesize
\begin{tabular}{@{}p{0.10\columnwidth}p{0.12\columnwidth}p{0.24\columnwidth}p{0.40\columnwidth}@{}}
\toprule
Dim. & Type & Unit & Computation \\
\midrule
$\BRreach$ & set & count & fixed point of $\pi$-enabled access from $x$ \\
$\BRflex$ & interval & kW, kVAr & $\sum_d \mathcal{C}_d(t;u_t,A_\pi(d))$ over reachable DERs \\
$\BRauth$ & duration & s & $\max(T_{\mathrm{cred}},T_{\mathrm{sess}},T_{\mathrm{cmd}})$ \\
$\BRphys$ & scalar & \emph{snapshot} p.u.-node; \emph{integrated} p.u.-node-min &
$\max_a J_V(F,u,a)$ at one operating point, or $\int \max_a J_V\,dt$ over the horizon \\
\bottomrule
\end{tabular}
\end{table}
